\chapter{The interview}
\label{ch:interview}

\section{What happens}

In every Abgabegespräch each group member gets a small task, comparable in difficulty to
the assignment tasks, and implements it live in the group's SWEB while the tutor watches
(official rules page). You also explain code --- usually a snippet of your own
implementation, often while you implement. If you are slow or cannot explain or
implement, you get an \emph{individual} deduction on the assignment's points. The tutor
stops you as soon as it is clear that you deserve the points: you do not have to finish a
polished solution, you have to show that you \emph{know where things happen and how to
change them correctly}.

Students' reports (see \file{a0/README.md}) agree on three things: the tasks are about
\textbf{syscalls}, the \textbf{page fault handler} and \textbf{page mappings}, threads and
\textbf{registers}; knowing your way around the code matters more than any single
algorithm; and the most common reason for deductions is \textbf{locking} --- missing or
wrong (chapter~\ref{ch:locking}).

\section{A strategy for the 30 minutes}

\begin{enumerate}
\item \textbf{Restate the task} in one sentence and ask about the edge cases you see
  (``what should happen if the pointer is invalid?''). Tutors expect questions.
\item \textbf{Say where it happens.} ``A segfault ends up in
  \code{PageFaultHandler::handlePageFault}, and for addresses without a segment in
  \code{Loader::loadPage} --- I'll hook in there.'' Chapter~\ref{ch:map} is the map for this step.
\item \textbf{List the pieces} before typing: syscall (6 places), per-process state (where?),
  page-fault branch, libc wrapper, test. Chapter~\ref{ch:recipes} has the checklists.
\item \textbf{Write the minimal correct version.} No refactoring of existing code, no
  generality nobody asked for.
\item \textbf{Locking pass.} For every piece of shared state you touched: which lock, where
  taken, where released, could anything run in between? Say it out loud.
\item \textbf{Test} with a tiny program in \file{userspace/tests/} --- and show the error path.
\end{enumerate}

\begin{infobox}[title=\textbf{Talking while coding}]
Practise narrating: what you look for, why this file, why this lock. It makes the tutor
see your understanding even when you have to look something up, and it keeps you from
silently typing yourself into a corner. The tasks' ``questions a tutor might ask'' are
there for this.
\end{infobox}

\section{How to use this course}

\begin{description}
\item[Examples] (\file{examples/}) --- one mechanism each, complete and tested; read the
  patch, apply it, run it. They are your reference when a task needs ``that thing with the
  ident mapping''.
\item[Tasks] (\file{a0/}, \file{a1/}, \file{a2/}) --- task cards in the tutors' format, a
  test program, hints behind fold-outs, questions. Time-box them.
\item[Solutions] (\file{solutions/}) --- patch plus \file{NOTES.md}. Read the \emph{Locking}
  and \emph{Pitfalls} sections even when your version works.
\item[Tools] --- you work in a separate training copy of SWEB, \file{repos/sweb-ag}, one git
  branch per task (created once with \code{task.sh setup}). \code{tools/task.sh} does the bookkeeping:
  \code{task.sh start a0-01} (fresh branch from the right base, tests copied in),
  \code{task.sh test} (build + run the tests without a window, with leak check),
  \code{task.sh run} / \code{kvm} / \code{debug} + \code{gdb} (exactly like \code{sweb-practise}),
  \code{task.sh reset a0-01} or \code{reset a1} (start over --- your attempt is kept in a
  \code{backup/} branch), \code{task.sh solution a0-01} and \code{task.sh diff a0-01}.
\end{description}
